\documentclass[12pt,a4paper]{article}

% Packages
\usepackage[utf8]{inputenc}
\usepackage[german]{babel}
\usepackage{geometry}
\usepackage{hyperref}
\usepackage{longtable}
\usepackage{graphicx}
\usepackage{parskip}

\geometry{a4paper, margin=1in}

% Title Page
\title{Pflichtenheft \\ \large Projektname}
\author{Dein Name}
\date{\today}

\begin{document}

\maketitle
\tableofcontents
\newpage

%--------------------------------------------------
\section{Einleitung}
\subsection{Ziel des Dokuments}
Dieses Pflichtenheft beschreibt die technische Umsetzung der im Lastenheft definierten Anforderungen.

\subsection{Projektübersicht}

Im aktuellen Prozess werden Liefertermine manuell durch Disponenten mit Kunden abgestimmt. 
Dieses Vorgehen ist zeitintensiv, fehleranfällig und nur begrenzt skalierbar.

Ziel des Projekts ist die Entwicklung eines digitalen, prototypischen Systems zur workflowbasierten Rückbestätigung von bereits durch die Disposition vorgeschlagenen Lieferterminen bzw. Lieferfenstern.

Das System übernimmt keine automatische Lieferterminplanung und keine automatische Tourenplanung. 
Die Planung und fachliche Entscheidung verbleiben bei der Disposition.

Das System unterstützt die Disposition bei der Anzeige vorgeschlagener Lieferoptionen, beim Versand von Rückbestätigungsanfragen sowie bei der strukturierten Verarbeitung von Kundenrückmeldungen.

Die Prozesslogik kann optional mithilfe von BPMN 2.0 modelliert und durch eine Workflow-Engine wie Camunda unterstützt werden.

%--------------------------------------------------
\section{Systemübersicht}
\subsection{Gesamtsystem}

Das Gesamtsystem dient der digitalen Rückbestätigung von Lieferterminen im DISPO-Umfeld. 
Es ersetzt nicht die Disposition und führt keine automatische Liefertermin- oder Tourenplanung durch.

Das System zeigt vorhandene Auftragsdaten, vorgeschlagene Lieferfenster und den aktuellen Rückbestätigungsstatus an. 
Über das System kann der Disponent eine Rückbestätigungsanfrage starten. 
Der Kunde kann das vorgeschlagene Lieferfenster bestätigen oder ablehnen.

Die Kundenreaktion wird gespeichert und der Status des Auftrags aktualisiert. 
Bei Ablehnung oder ausbleibender Rückmeldung wird der Auftrag als problematischer Fall markiert und dem Disponenten zur manuellen Weiterbearbeitung angezeigt.

Optional kann eine Workflow-Komponente eingesetzt werden, um den Rückbestätigungsprozess technisch zu steuern.

\subsection{Systemkontext}

Das System wird als eigenständige Anwendung im Umfeld eines bestehenden DISPO-Systems betrieben.

Auftragsdaten und Lieferinformationen werden aus dem DISPO-Umfeld übernommen und im Rückbestätigungsprozess verarbeitet.

Das System ersetzt keine bestehende Dispositionssoftware und übernimmt keine automatische Liefertermin- oder Tourenplanung.

\subsection{Systemarchitektur}
\begin{itemize}
    \item Frontend (Web UI für Disponenten und ggf. Kundenseite)
    \item Backend (Geschäftslogik und REST-API)
    \item Datenbank
    \item E-Mail-Kommunikationsschnittstelle
    \item Optional: Workflow Engine, z.B. Camunda
\end{itemize}

%--------------------------------------------------
\section{Prozessbeschreibung}
\subsection{Soll-Prozess}

Der Soll-Prozess beschreibt die digitale Rückbestätigung eines bereits durch die Disposition vorgeschlagenen Lieferfensters.

Der Prozess umfasst folgende Schritte:

\begin{enumerate}
    \item Der Disponent wählt einen bestehenden Auftrag mit vorgeschlagenem Lieferfenster aus.
    \item Das System zeigt Auftragsdaten, Lieferfenster und Tourinformationen an.
    \item Der Disponent startet eine Rückbestätigungsanfrage.
    \item Das System versendet die Anfrage primär per E-Mail.
    \item Der Kunde bestätigt oder lehnt das vorgeschlagene Lieferfenster ab.
    \item Das System speichert die Kundenreaktion.
    \item Der Rückbestätigungsstatus des Auftrags wird aktualisiert.
    \item Abgelehnte Aufträge oder Aufträge ohne Rückmeldung werden dem Disponenten als problematische Fälle angezeigt.
\end{enumerate}

Eine automatische Planung neuer Lieferfenster oder Touren ist nicht Bestandteil des Systems.

%--------------------------------------------------
\section{Datenmodell}
\subsection{Übersicht}

Beschreibung der zentralen Datenobjekte:

\begin{itemize}
    \item Auftrag
    \item Kunde
    \item Lieferfenster
    \item Tourinformation
    \item Rückbestätigung
    \item Status
\end{itemize}

%--------------------------------------------------
\section{Funktionsbeschreibung}

\subsection{Anzeige von Aufträgen}

Das System zeigt bestehende Aufträge mit geplantem Lieferdatum, geplantem Lieferfenster und zugeordneter Tourinformation an.

\subsection{Start einer Rückbestätigungsanfrage}

Der Disponent kann für einen ausgewählten Auftrag manuell eine Rückbestätigungsanfrage starten.

\subsection{Benachrichtigung des Kunden}

Das System versendet die Rückbestätigungsanfrage primär per E-Mail. 
Weitere Kommunikationskanäle können optional ergänzt werden.

\subsection{Kundenreaktion}

Der Kunde kann das vorgeschlagene Lieferfenster bestätigen oder ablehnen.

\subsection{Statusverarbeitung}

Das System speichert die Kundenantwort und aktualisiert den Rückbestätigungsstatus des Auftrags.

\subsection{Behandlung ausbleibender Rückmeldungen}

Wenn innerhalb einer definierten Frist keine gültige Antwort eingeht, setzt das System den Status auf ``keine Rückmeldung''.

\subsection{Anzeige problematischer Fälle}

Abgelehnte Aufträge und Aufträge ohne Rückmeldung werden dem Disponenten als problematische Fälle angezeigt.

\subsection{Manuelle Weiterbearbeitung}

Der Disponent kann problematische Fälle außerhalb des automatisierten Rückbestätigungsprozesses manuell weiterbearbeiten.

%--------------------------------------------------
\section{Technische Funktionsanforderungen}

\begin{longtable}{|p{1.5cm}|p{4cm}|p{8cm}|}
\hline
\textbf{ID} & \textbf{Funktion} & \textbf{Technische Umsetzung} \\
\hline

PF1 & Anzeige von Aufträgen &
Das Frontend ruft Auftragsdaten über eine REST-API ab und zeigt Lieferdatum, Lieferfenster, Tourinformation und Rückbestätigungsstatus an. \\
\hline

PF2 & Start einer Rückbestätigungsanfrage &
Der Disponent kann über das Frontend eine Rückbestätigungsanfrage starten. Das Backend verarbeitet die Anfrage und erstellt eine eindeutige Rückbestätigungs-ID. \\
\hline

PF3 & Versand der Rückbestätigungsanfrage &
Das Backend versendet automatisiert eine E-Mail mit einem Weblink zur Rückmeldung an den Kunden. Weitere Kommunikationskanäle können optional ergänzt werden. \\
\hline

PF4 & Verarbeitung von Kundenreaktionen &
Kundenreaktionen werden über einen Weblink entgegengenommen, durch das Backend verarbeitet und in der Datenbank gespeichert. \\
\hline

PF5 & Statusverwaltung &
Das System aktualisiert den Rückbestätigungsstatus abhängig von der Kundenreaktion automatisch in der Datenbank. \\
\hline

PF6 & Behandlung ausbleibender Rückmeldungen &
Wenn innerhalb der definierten Frist keine gültige Antwort eingeht, setzt das Backend den Status automatisch auf ``keine Rückmeldung''. \\
\hline

PF7 & Anzeige problematischer Fälle &
Aufträge mit Ablehnung oder fehlender Rückmeldung werden im Frontend als problematische Fälle dargestellt. \\
\hline

PF8 & Verhinderung mehrfacher Antworten &
Das System stellt sicher, dass pro Rückbestätigungsanfrage nur eine gültige Antwort statuswirksam gespeichert wird. \\
\hline

PF9 & Workflow-Unterstützung (optional) &
Optional kann eine Workflow-Engine wie Camunda verwendet werden, um den Rückbestätigungsprozess technisch zu steuern. \\
\hline

\end{longtable}

%--------------------------------------------------
\section{Geschäftslogik}
\begin{itemize}
    \item Ein Auftrag kann eine Rückbestätigungsanfrage erhalten.
    \item Pro Anfrage darf nur eine gültige Kundenantwort statuswirksam gespeichert werden.
    \item Bei Bestätigung wird der Status auf ``bestätigt'' gesetzt.
    \item Bei Ablehnung wird der Status auf ``abgelehnt'' gesetzt.
    \item Bei ausbleibender Antwort innerhalb der Frist wird der Status auf ``keine Rückmeldung'' gesetzt.
    \item Abgelehnte Aufträge und Aufträge ohne Rückmeldung werden als problematische Fälle markiert.
    \item Das System führt keine automatische Liefertermin- oder Tourenplanung durch.
\end{itemize}

%--------------------------------------------------
\section{Schnittstellen}
\subsection{Interne Schnittstellen}
\begin{itemize}
    \item Frontend – Backend
    \item Backend – Workflow Engine
    \item Workflow Engine – Datenbank
\end{itemize}

\subsection{Externe Schnittstellen}
\begin{itemize}
    \item E-Mail-Service für den Versand von Rückbestätigungsanfragen
    \item Weblink für die Kundenrückmeldung
    \item Optional: weitere Kommunikationskanäle, z.B. WhatsApp
\end{itemize}

%--------------------------------------------------
\section{Testkonzept}
\subsection{Teststrategie}
Unit-, Integrations- und Prozesstests.

\subsection{Testfälle}
\begin{itemize}
    \item Kunde bestätigt Lieferfenster: Status wird auf ``bestätigt'' gesetzt.
    \item Kunde lehnt Lieferfenster ab: Status wird auf ``abgelehnt'' gesetzt.
    \item Kunde antwortet nicht innerhalb der Frist: Status wird auf ``keine Rückmeldung'' gesetzt.
    \item Kunde antwortet mehrfach: Nur eine gültige Antwort wird statuswirksam gespeichert.
    \item Abgelehnte Aufträge werden als problematische Fälle angezeigt.
\end{itemize}

\end{document}